\documentclass[11pt,a4paper]{article}

\usepackage{iftex}
\ifPDFTeX
  \usepackage[utf8]{inputenc}
  \usepackage[T1]{fontenc}
  \usepackage{lmodern}
\fi
\usepackage{amsmath,amssymb,amsthm}
\usepackage[margin=25mm]{geometry}
\usepackage{microtype}
\usepackage[hidelinks]{hyperref}

\newtheorem{theorem}{Theorem}
\theoremstyle{definition}
\newtheorem{definition}[theorem]{Definition}
\theoremstyle{remark}
\newtheorem{remark}[theorem]{Remark}

\title{\textbf{The Theorem of One Bit and Two Friends}}
\author{Jan Mikul\'ik}
\date{27 November 2026}

\begin{document}
\maketitle

\begin{abstract}
Two friends start from the same history and ask the same first question; the
prescribed answer is a fair bit. Each may then choose a second question from the
first answer; the prescribed answer to every further question is \(0\). One fair
bit is an exact causal seed of entropy \(1\), and every exact seed has entropy
at least \(1\). A coupling with the correct marginals need not be the joint of a
shared causal decoder. Lean checks the identities in
\texttt{CausalSeed/OneBitFriends.lean}.
\end{abstract}

\section{The family}

A deterministic adaptive policy \(\pi\) maps the first reply \(b\in\{0,1\}\) to
a second question \(\pi(b)\) in a set \(A\) of actions. The prescribed
transcript of \(\pi\) on first reply \(b\) is \((b,\pi(b),0)\). The prescribed
law places mass \(1/2\) on each of the two transcripts
\((0,\pi(0),0)\) and \((1,\pi(1),0)\).

A finite seed \(\mu\) on a type \(\alpha\), together with a first-reply map
\(\varphi:\alpha\to\{0,1\}\), induces, under a policy \(\pi\), the transcript
\((\varphi(s),\pi(\varphi(s)),0)\) on each atom \(s\). The pair \((\mu,\varphi)\)
is \emph{exact} on \(A\) when this law equals the prescribed law of every
policy. Later replies are the constant \(0\).

\section{Sufficiency}

\begin{theorem}
The fair bit \(S\) with the shared decoder that returns \(S\) on the first
question and \(0\) thereafter is an exact causal seed. For every policy
\(\pi\), the transcript \(X_\pi=(S,\pi(S),0)\) has the prescribed law. The seed
has Shannon entropy \(H(S)=1\). Every realized transcript of every policy has
surprisal \(1\), so the expected surprisal of every policy is \(1\), and the
envelope is \(L=1\).
\end{theorem}

\section{Necessity}

\begin{theorem}
Let \(A\) be nonempty. If a finite nonnegative seed \(\mu\) and a first-reply
map \(\varphi\) are exact on \(A\), then the first reply \(B'=\varphi(S')\) is a
fair bit. The entropy of a function is at most the entropy of the seed, so
\(H(S')\ge H(B')=1\). Together with the previous theorem, the causal cost is
\(C=1\).
\end{theorem}

\begin{remark}
The hypothesis that \(A\) is nonempty is used only to produce a policy, so that
exactness can be evaluated and the first-reply law identified with the fair bit.
The lower bound \(H(S')\ge 1\) itself holds for every finite nonnegative seed
whose first reply is a fair bit.
\end{remark}

The quantities \(L=1\) and \(C=1\) are the expected surprisal of every policy
and the exact-seed entropy characterization of this family. The equality
\(L(T)=C(T)=1\) on the two-friend tree is already proved in
\texttt{CausalSeed/ManyFriends.lean}.

\section{Correct marginals}

\begin{theorem}
Let one friend ask a constant action \(a\) and the other a constant action
\(b\neq a\). The anti-coupling that records \((U,a,0)\) for the first friend and
\((1-U,b,0)\) for the second, with \(U\) a fair bit, has both prescribed
marginals. No shared causal decoder induces this joint: a shared decoder returns
the same first reply, at the same history, on every seed atom, whereas the
anti-coupling places mass \(1/2\) on a pair of opposite first replies.
\end{theorem}

\begin{remark}
The hypothesis \(b\neq a\) is part of the stated scene of two distinct second
questions. The contradiction uses only that the anti-coupling has opposite first
replies.
\end{remark}

Not formalized: a general tree, and the greedy overhead.

\end{document}
